\section{Introduction}
A battery-free device must turn an intermittent supply into a particular
service: capture an external event, preserve the relevant state, and deliver
the required information at the required time. Having enough energy in total
does not specify whether a rail becomes ready before a deadline, whether a
transmit pulse causes a reset, or whether the event is still pending when the
device recovers. The interface presented to a designer determines which of
these relations can be evaluated before a legal change is made.

The distinction matters even when the current service succeeds. Two worlds
can produce the same present voltage and event traces while having different
unexcited storage clamps. Reducing capacitance activates those clamps and
separates their outcomes. Two pending-event memories can likewise behave
identically until a reset is introduced. These examples do not imply that
static descriptions are intrinsically inadequate. A static description of the
relevant laws, parameters, initial state and legal intervention domain can
already answer the prospective question~\cite{structural,mechanisms}.

An observation interface offers another route: acquire a source, select or
construct a relation, qualify its domain and units, and feed the result into
a design consumer. That route has an operational burden of its own. A useful
equation with an invalid path can fail admission; a correctly computed Boolean
can fail to match a numeric consumer branch. A qualified observation can also
leave the final design unchanged, or support a design that complete static
information already permits. Source acquisition, typed consumption, service
success and incremental benefit therefore require distinct evidence~\cite{contract,original,tables}.

We study these distinctions with sealed, executable digital BF cases. The
first part states the information requirement and instantiates two structural
counterexamples. The second implements the acquisition-to-consumer route and
compares it against average-energy, rich-static, complete-input and executable
interfaces under matched budgets. The main formal comparison has four fixed
tasks and three repetitions per condition. Its strongest negative result is
central: rich-static information already obtains 12/12 complete-service passes,
whereas live observation completes seven chains and leaves five groups
incomplete. The equal-information serialization comparison gives the same
12/12 service and 65 requests in both arms~\cite{comparison,tables}.

A secondary vibration-event study tests a different service boundary. Correct
event matching and energy accounting do not ensure correct encoded event
times. Two designs fail one-second timestamp tolerances even though their
messages arrive promptly. We retain these submitted failures separately from
four token-censored episodes with no design submission~\cite{final}.

The contribution is a connected account of what information must distinguish,
how an observation must reach a real consumer, and where successful digital
checks stop supporting the claim. It includes the theoretical condition,
mechanism witnesses, strong static controls, actual failed and completed
observation chains, and portable saved-data reconstruction. It supplies no
claim of priority over general sufficient-statistic or abstraction theory,
no advantage over the successful complete static controls, and no hardware
qualification. The study concerns public development apparatus rather than
an official benchmark or a population of unseen physical devices.
